\section{The Road to the true West}

We will return to the Number 8 \& Iamblichus, very soon.

The Inklings are well-known to Christians, and greatly lauded; unfortunately, very few are interested in the more profound portions of any of their work.

Gornahoor has already pointed out\footnote{\url{http://www.gornahoor.net/?p=3545}} that CS Lewis ultimately held a reductive view of the myths which he so delighted in; however, \emph{it may have been somewhat different for Tolkien.} They had made a pact together: Lewis was to write a novel about space-travel (in which he suggested that the celestial elementals could be real -- did he mean this?), and Tolkien would write a story about time-travel. Tolkien became sidetracked with Lord of the Rings, but ``The Lost Road'' was a story which he tried to weave into the Mythos of LOTR.

It begins with a father and son, both Englishmen; the son is hearing voices through the natural elements, or in dreams, which spell out the syllables of what he thinks is a language. Alboin is the son's given Christian name. Alboin is dark-complected, and going off to the university -- his mother is dead, and he is an only child. He is staying with his ailing father just prior to going off to school, and is introspecting:

\begin{quotex}
``Dark Alboin. I wonder if there is any Latin in me. Not much, I think. I love the western shores, and the real sea -- it is quite different from the Mediterranean, even in stories. I wish there was no other side to it. There are darkhaired people who are not Latins. Are the Portugese Latins? What is Latin? I wonder what kind of people lived in Portugal and Spain and Ireland and Britain in the old days, very old days, before the Romans and the Carthiginians. Before anybody else. I wonder what the man thought who was the first to see the western sea.''
\end{quotex}


His father argues the next day:

\begin{quotex}
``Of course, I am not a philologist, but I could never see that there was much evidence in favour of ascribing language-changes to a substratum. Though I suppose underlying ingredients do have an influence, though it is not easy to define, on the final mixture in the case of peoples taken as a whole, different national talents and temperaments, and that sort of thing. But races, and cultures, are different from languages.''
\end{quotex}


Alboin argues with his father about prehistory -- where his father thinks that pre-culture and pre-history would be beastly and animal-like, Alboin is inclined rather to think that something better and greater inhabited the time of the past along the coasts of Cornwall. They end up discussing the ``secret language'' that Alboin is deciphering from intuitions: \emph{lomelinde}, for example, he is pretty sure means nightingale.

One of the Eressean Elf-latin phrases he had translated was this:

\emph{Westra lage wegas rehtas, nu isti sa wraithas}.


\textbf{``A straight road lay westward, now it is bent''.}


Alboin's father dies over the course of early term, and Alboin is saddened because they had not been as close immediately prior to his death as they had when he was younger, and had just started the translations. This turns out to be reflective of a father-son relationship that was occuring in the hyper-reality of the world which Alboin is dreaming about.

\begin{quotex}
Thus said Aelfwine\footnote{\url{http://www.elvish.org/gwaith/aelfwine.htm}} the far-traveled -- ``There is many a thing in the West regions unknown to men, marvels and strange beings, a land fair and lovely, the home land of the Elves, and the bliss of the Gods. Little doth any man know what longing is his whom old age cutteth off from return.''
\end{quotex}


Before his father dies, he tells him -- you can't go back.

Alboin however, realizes that his inner being is differently constituted:

\begin{quotex}
``He felt he could say that his most permanent mood, though often overlaid or suppressed, had been since childhood, the desire to go back. To walk in Time, perhaps, as men walk on long roads; or to survey it, as men may see the world from a mountain, or the earth as a living map beneath an airship.''
\end{quotex}


Eventually, he deciphers a long passage:

\begin{quotex}
And Sauron came to Numenor, fell under Shadow, war-made on powers -- sea fell into the Chasm and Numenor down-fell. The road straight went west-ward, all roads were now bent, and the death-shadow is heavy (to us).
\end{quotex}


As the story unfolds in the hyper-reality, it becomes clearer that at one point, a straight road into the West passed gradually upwards into the lower heavens through the lower waters, and that men once walked this road to the abode of the gods; after the coming of Sauron, and the Fall, when Atlantis was overthrown, the world became round, and the road into the West was bent, so that only after death could men walk it. This is part of the story of Numenor, which occurs before the Lord of the Rings tale. Numenor and Atlantis are linked.

Alboin has had a son, named Audoin, who is also hearing visions and dreaming dreams. He is too young to know that in the past, a figure had appeared in a dream to Alboin, calling himself Elendil, the elf-friend, who offers him his desire : to go back.

\begin{quotex}
``That cannot be, it is against the law.''

``It is against the rule -- laws are commands upon the will and are binding. Rules are conditions; they may have exceptions. They may be strict,yet they are the means, not the ends, of government. There are exceptions, for there is that which governs, and is beyond the rules. Behold, it is by the chinks in the wall that the light comes through, whereby men become aware of the light and therein perceive the wall and how it stands. The world is not a machine that makes other machines, after the fashion of Sauron. To each under the rule, some unique fate is given, and he is excepted from that which is a rule to others.''
\end{quotex}


Alboin answers, then, that yes, he will have his desire. However, his very young son, Audoin, is drawn into the plot.\footnote{\url{http://en.wikipedia.org/wiki/Audoin}}

\begin{quotex}
Then Alboin seemed to fall into a dark and a silence, deep and absolute. It was as if he had left the world completely, where all silence is on the edge of sounds, and filled with echoes, and where all rest is but repose upon some greater motion. He had left the world and gone out. He was silent and at rest: a point.
\end{quotex}


There is a father-son pair in Numenor, long ago, who are experiencing a similar struggle together -- the father is trying to get the son to disavow allegiance to Sauron, who is now effectively controlling their king. The paternal bond prevails, the son, against his will, disavows Sauron. This parallels Audoin being dragged into Alboin's choice in spite of his free will, due to the bargain and the blood relationship.

``There is storm over Numenor.''

The fragments are not all linked by Tolkien, so this is not a ``Story'' per se. However, he had interesting thoughts and designs for this one.

\begin{quotex}
``The Father made the world for elves and mortals, and he giveth it into the hands of the Lords, who are in the West. There is Iluvatar, the One, and then there are the Powers, of whom the eldest in the thought of Iluvatar was Alkar the Radiant; and there are the firstborn of Earth, the Eldar, who perish not while the world lasts. And then there are also the afterborn, mortal men, who are the children of Iluvatar, and yet under the rule of the Lords. Iluvatar designed the world, and revealed this design to the Powers; and of these, some of these he set to be Valar, Lords of the World, and governors of the things that are therewith. But Alkar, who had journeyed alone in the Void before the World, seeking to be free, desired the World to be a kingdom unto himself. Therefore, he descended into it like falling Fire; and he made war upon the Lords, his brethren. But they established their mansions in the West, in Valinor, and they shut him out, and they gave battle to him in the North, and they bound him, and the World had peace, and became exceeding fair... all things have an end in this world, for the world itself is bounded, that it may not be void. But death is not decreed by the Lords: it is the gift of the One, and a gift which in the wearing of Time even the Lords of the West shall envy.''
\end{quotex}


Elendil teaches his son that though wisdom is lost in Numenor with Sauron, we can have the wisdom to know there is no escape, except to a worse fate: ``the old songs are altered or forgotten; twisted into other meanings''. Sauron offers Empire, Machines, and Conquest to the West, but Elendil persuades Herendil to join him in disobedience to Sauron (and in obedience to the King, although his mind is twisted): ``Even to dispraise Sauron is held to be rebellious.''

Does this not describe our own era?

Sauron cannot build ships that sail the road upward into the West, although the Numenoreans tried to build airships previous to this, which never reached the mansions of Valinor. All the Empire can do is to build engines of war and conquest that are limited to the horizontal level.

The world, which was once a straight road up into the heavens, a vertical ascent on a level incline, has become bounded: a fallen orb in which the vertical is chained to the horizontal, and the true road West is closed. Nevertheless, this is not Fate or a Law, but a Rule, which has exceptions. Certain exceptional men, such as Aelfwine or Alboin, can have expriences of Eressea and the Western Isles, but this is the exception now, not the rule. Some can escape, because, as Tolkien intimates, the true road to the West is now hidden in the post-mortem states, and is therefore still existent, even in the ``round world'' of waking men.

Tolkien goes on to link his stories of the West to the voyages of St. Brenden\footnote{\url{http://en.wikipedia.org/wiki/Brendan\#The_Voyage_of_Saint_Brendan}}.

I hope one can get a small taste of how a master storyteller, faithful Catholic, and man of letters was able to go further than someone like Lewis in intuiting how the old myths did not breathe lies through silver\footnote{\url{http://en.wikipedia.org/wiki/Mythopoeia_\%28poem\%29}}, but actually were ``true myths''; wherever Tolkien is, I am certain Valinor and Middle Earth exist, even if he has passed beyond them. And if they exist in the heavens, then there is also a road into the true West, however little traveled, that is accessible still. Tolkien's myths hold great potential teaching power, when harmonized and explained within the Tradition.


\flrightit{Posted on 2013-03-15 by Logres}

\begin{center}* * *\end{center}

\begin{footnotesize}
\begin{sffamily}
\texttt{Scardanelli on 2013-03-16 at 10:01 said:}

Great job... this was quite refreshing and inspirational. im glad you linked to the Mythopoeia poem. I only discovered a few days ago and it was somewhat of a revelation reading it. Coincidently, I've just begun reading B.M.'s Gnosis series and was reading some of your previous posts on it. I've only just begun the first book, but it seems like there are some similarities here between the concept of surpassing the ``general law'' and Tolkien's concept of rules as opposed to laws. Anyway, thanks for this!

\hfill

\texttt{Thorsten on 2013-03-16 at 15:27 said:}

I need to go reread Tolkien's works (especially such lesser-known ones as The Book of Lost Tales with which the story at hand is evidently connected) having since learned about the metaphysical significance underlying his mythology. As a philologist myself I can't help but point out that the ``Elf Latin'' quotation,

``Westra lage wegas rehtas, nu isti sa wraithas.''

is what appears to be Tolkien's rendering of Proto-Old English or some stage intermediate between Proto-Germanic and Old English. Given that Tolkien's use of language both historical and constructed (or in this case reconstructed) often has multiple layers of symbolism (E.g. Rohirric language and culture drawing so heavily from that of the Anglo-Saxons) and often alludes to historical or mythical motifs (E.G. Eärendil -- ?arendel `luminous wanderer' of Germanic Myth:\url{http://en.wikipedia.org/wiki/Earendel}) I would be interested in knowing more about the context of this quotation, because the use of this phase of the English language not used elsewhere by Tolkien to my knowledge.

All the best.

\hfill

\texttt{Logres on 2013-03-16 at 21:23 said:}

At the moment, I've misplaced my only copy, so I'll try to post the context of the quote, Thorsten, later, but it shows up in several places in the opening of The Lost Road.

Scardanelli, very insightful: yes, I put that quote in because it is a beautiful exposition both of the unique destinies of the individual, and also because it draws a very subtle and necessary distinction between Rule and Law. General Law attempts to keep ``everyone in their place'' : the true individual will ``escape'' General Law by proper and lawful efforts (to which Mouravieff gives directions) which place them on the path of destiny, which is the exception to the rule because it fits the emerging hero. A moment of reflection will show that this happens all the time on the non-supernatural level anyway, and is designed to lead us to this great realization.

\hfill

\texttt{Logres on 2013-03-16 at 21:38 said:}

Alboin translates that passage along with the Aelfwine quotation for his father, just before the old man goes into his final decline. So it is a last scene of the ``good old days'' (discussing philosophy, etc.) between father and son. And thank you for the links, sir.

\hfill

\texttt{Thorsten on 2013-03-17 at 10:13 said:}

Thank you for the further info, Logres. I'm always happy to provide input regarding things of this sort. OE and Germanic philology is one of my areas of focus on account of ancestry and appreciation of the insight it provides into older Indo-European spirituality, especially regarding the role of the warrior, something Tolkien understood so thoroughly.

\hfill

\texttt{J. Alexander Maximilian on 2013-03-17 at 14:05 said:}

That was excellent!

Tolkien really had a great intuition and imagination. The vision I got just from the small anecdotes, incredible!

The synthetic, technological empire by which Sauron imposes his rule (as you stated is the state of the world in our time). In contrast to the organic, initiatory, and pure existence the men of the west exist in is superb. The visions of Alboin (whose name seems to be reference to Albion), which arise from his inner souls memories (intuition), and his struggle to understand them in reference to the prevailing thinking (his Father) about myth and legend, within the standard contemporary religious, and secular views. As opposed to what he intuitively knows to be the case.

Tolkien's presenting of myths and legends of the ancient past through a devoted Catholic worldview as not being mere folk tales and, or pagan devilry, that should be either enjoyed as fiction- or completely disregarded as the latter. But as living within the spiritual memory of us in the now, and being purified by the blood of Christ point to a higher path, and truth that has been lost to most, but can be recalled by those who have the intuition, and understanding to do so. Both of which are gifts of the Holy Spirit, but thankfully can be strengthened by contemplation, and Sacramental life.

I thank you for this!

\hfill

\texttt{Logres on 2013-03-17 at 22:07 said:}

I was shocked, also, when I read the Lost Road. Pleasantly so.

\hfill

\texttt{J. Alexander Maximilian on 2013-03-18 at 12:24 said:}

Indeed!

\hfill

\texttt{Logres on 2013-03-18 at 13:04 said:}

I would like to see a traditional artist of letters take the idea of The Lost Road, and integrate it into the Tradition as purest art. It would be like reforging an ancient sword.

\hfill

\texttt{J. Alexander Maximilian on 2013-03-18 at 13:29 said:}

Hmmm... Well then, I may have to make an attempt.This would be a most excellent task, I do write a bit of poetry, mostly prose.

Creating a vision of the Great struggle between cthonhic techo/ materialistic forces and the Solar organic initiatory forces, played out somehow in the inner, and exterior world of a character existing in the contemporary world. Who must overcome himself, and the society at large, with some sort of external aid that exists in the greater present, which is now, and simultaneously all times. And it could incorporate the ideas of great Traditionalist thinkers, and ancient Spiritual wisdom. The trick is to take the concept without being a complete mimicry.

\hfill

\texttt{Scardanelli on 2013-03-18 at 13:42 said:}

I haven't had much time to explore this concept, but lately I've thought that it would be interesting to use traditional forms of verse but in a modern context... For instance epic poetry written in an old English alliterative verse form but dealing with the true heroic in the modern world. It could serve as a type of bridge... 

\hfill

\texttt{scardanelli on 2013-03-18 at 13:54 said:}

If enough people involved themselves in this effort, an entire descriptive language could be developed that would help us identify the truth of what the modern world really is, somewhat along the lines of Old English kennings used in reference to themes of Old English /Norse poetry. Tolkien's use of the ``Iron Crown'' in Mythopoeia is a great example.

\hfill

\texttt{J. Alexander Maximilian on 2013-03-18 at 13:56 said:}

Yes. The heroic in action in this darkest of times. For now the hero must stand against the entire current of wrong thinking so indulged in, and enforced. The counter- initiatory motion of ``progress'' and the adharmic thought called ``liberation'' in its various forms. The monsters he would be up against are feminism, homosexual ``rights'', etc. And these could be personified with great villainous beasts that maintain their existence by the psychic vampirism of their hosts, that being the special interest group that it feeds upon, and these are all merely pawns used by the ultimate dark lord whose goal is global empire, and his magick is ``sin''-ema, and television, etc.

\hfill

\texttt{J. Alexander Maximilian on 2013-03-18 at 13:57 said:}

A collaborative effort would be good.

\hfill

\texttt{scardanelli on 2013-03-18 at 14:58 said:}

I think that it would be useful to remember Tolkien's distaste for allegory here. His writings are true creation, or subcreation in his words, and do not necessarily symbolize or personify anything. His main concern was to creat a truly moving work of art. To me, it seems somewhat akin to Tomberg's description of sacred magic... to create something and project it into the world and which consists of a unity between one's will and the divine will.

\hfill

\texttt{J. Alexander Maximilian on 2013-03-18 at 15:07 said:}

Yes. However these anti- traditional movements are beasts that exist on the energy of the people involved in them, and are ultimately tied to the agenda of mono-cultural globalism, which are run by black magicians, who use spells through the media. The media is black magick, that conforms most people to the agenda that it promotes. And it would be a useful way to show the inherent lower nature of these ideas that are promoted by these agents, and activists.

\hfill

\texttt{Thorsten on 2013-03-18 at 17:13 said:}

True enough. Tolkien did profess distaste for allegory on more than one occasion, particularly when some suggested outright, and simple-minded equivalences with historical or contemporary developments (e.g. between Saruman and Hitler, or the orcs and German troops). It goes without saying that Tolkien was operating on an much more profound level than something so cheap and superficial as that. However the point I make is that Tolkien did (as can easily be verified with many many examples) use real world languages, proper names from real world history, etc as stand ins for the ``actual'' languages, proper names etc. of Middle Earth. This was done because he reckoned it would carry greater evocative power for his readers than would the use of the constructed ones. Tolkien said plainly that the use of Old English for the language of the Rohirrim was a surrogate for the real language of Rohirric. Old English and Old English names were felt to carry much of the same character, and to express the spirit or inner nature of the Rohirrim. Analagous equivalences abound, such as the use of Hobbit for khu-dukan etc.

So the use of Proto-Germanic in that phrase has some significance. It may be a subtle one, perhaps even relatively in significant in the grand scheme of the narrative, but nonetheless no accident.

I didn't mean to dwell on this aspect more than is called for (for now anyway) but being a lover of language myself my mind immediately homes in on these things and am aware of how important languages were to Tolkien for their ability to suggest the world-view of a particular ``race'' and the like.

\hfill

\texttt{Alexander Maximilian on 2013-03-18 at 18:22 said:}

oops. I thought you were replying to me. Sometimes its hard to tell exactly what's going on with the phone view.

\hfill

\texttt{Logres on 2013-03-24 at 06:31 said:}

I meant to link to this wonderful post:


\url{http://failedhermit.blogspot.com/2012/12/catholic-poem-in-time-of-war.html}

\hfill

\end{sffamily}
\end{footnotesize}

